\chapter{Difference Becomes Countable}

\section{Barely more than telling apart}

The ground floor left the machine holding decided differences and nothing else --- positions in an alphabet, each
tanged apart from the rest, none of them yet a number. This chapter takes the first step off that floor, and the step
is smaller than it looks. To count is not to invent a new faculty; it is to \emph{funge} --- to pool the marks the
machine has already tanged apart into a single bin, and read off how many the bin holds. The count is barely more than
telling apart. It falls out of the distinction the moment the machine is allowed to gather.

Watch the gathering exactly. The machine has decided some marks the same and others different; funging pools every mark
that decides the same way into one bin and keeps the bins that decide differently apart. That is the whole operation:
a decided sameness pooled, a decided difference kept. And once the marks are pooled, a number is simply the size of a
bin --- the cardinality of a class the machine assembled by nothing but its two verbs. Counting is admissible here
precisely because it asks for no power the ground floor did not already have: the machine tanged the differences, it
funges the sames, and the count is the tally of the funge. Nothing continuous has entered; nothing has been assumed.
The number is a finite bin-size, decided all the way down.

It is worth being exact about what makes a count \emph{admissible} in this construction, because the word is doing
work. A count is admissible when every mark in the bin was decided in, one at a time, by a procedure that halted --- and
inadmissible the moment it would require deciding infinitely many marks at once, or pooling by a characteristic the
machine cannot settle. The machine counts only what it can enumerate: finite bins, filled by halting decisions,
tallied by a funge that closes. A cardinality it could reach only by an oracle --- by helping itself to a completed
infinity it never enumerated --- is not a count this machine performs. Countable, for this machine, means exactly
\emph{reachable one halting decision at a time}, and the count that falls out of the distinction is always of that
humble, admissible kind.

In the two verbs the count is the funge's first fruit. The tange gave the marks apart; the funge pools the ones that
decide alike into a bin; and the number is the bin read off. Where the ground floor could only tange --- select this
mark against that --- this chapter adds the funge and gets counting for free, because a count \emph{is} a funge
tallied. The machine did not learn to count so much as it learned to gather what it had already told apart, and discover
that the gathering has a size.

\section{A number is a finite condition}

A bin-size is a whole number, and whole numbers are the easy case. The chapter's real work is the hard one: what does
this machine mean by a number that is \emph{not} whole --- a value between the marks, the kind the later chapters will
need when a bound falls between two counts? The answer is the deepest idea in the chapter, and it governs the whole
book. A number, to this machine, is not a completed object sitting out in some continuum. It is a \emph{finite
condition} --- a specification the machine can write down in full, that pins the value to an interval and no finer.

Take the construction Dedekind~\cite{dedekind1872} gave, and read it as the machine reads it. A real number is fixed by
a \emph{cut}: a rule that sorts every rational into those below it and those above. The machine cannot hold the whole
cut at once --- that would be an infinity of decisions completed --- but it does not need to. It holds a \emph{finite}
condition: a lower rational known to fall below and an upper known to fall above, a bracket the machine can actually
write down. Each further decision the machine can afford tightens the bracket --- narrows the gap between the walls ---
without ever closing it to a point. The value is approached through a tower of finite conditions, each one a thing the
machine genuinely computed, and the number \emph{is} that tower, not any limit lying past its end. Cantor's~\cite{cantor1872} account of the reals as sequences of rationals converging on a value says the same from the sequence side; this
machine keeps only the finite front of the sequence, the part it has actually reached.

The point to hold is that the limit is never completed, and this is a feature, not a shortfall. A machine that claimed
to hold the finished real --- every digit of it, the cut sorted to the last rational --- would be claiming to have run
an infinity of decisions to a halt, which is exactly the oracle the preface refused. This machine claims less and can
deliver it: a finite condition, a bracket around the value, tightened as far as its resolution allows and no further.
The value is real, and the machine's grip on it is honest, because the grip is finite. When the far chapters need a
number that falls between the counts, this is the number they will get: not a point fetched from the continuum, but a
tower of finite conditions, approached and never finished.

There is an older shape lurking under this, and it is worth naming because the whole construction will turn out to wear
it. Cohen~\cite{cohen1963}, proving a question of set theory independent, approached an ungraspable object through an
ascending tower of \emph{finite conditions} --- each one a thing you could write down, converging on a generic limit
that no single condition ever reached. That is precisely the shape of a number here: a generic value, approached
through finite conditions, never captured by any one of them. The machine will meet this shape again at the very end,
when it hands back a bracket instead of a point; the bracket is a finite condition, and the value it brackets is the
generic limit the finite conditions approach and honestly never complete.

In the two verbs a number is a funge held open. Each finite condition funges together every value consistent with the
bracket so far --- pools, under one interval, all the reals the machine cannot yet decide apart --- and each tightening
tanges the interval finer, splitting off the values a new decision rules out. A whole number is the case where the
funge closes on a single bin; a real is the case where it never quite does, and the machine keeps the open bin as its
honest answer. Number, in this machine, is the funge of the indistinguishable held to a finite width --- which is the
one shape the whole book is built to defend.

\section{The first loop}

The machine can tell apart and it can count; the last thing this chapter needs is the fact that these survive being
\emph{re-expressed}. A number the machine builds must remain the same number when the machine encodes it, ships it
through its own representation, and reads it back. This is the first loop, and it is the humble ancestor of every loop
the book will close, up to the great one at the end. The test is simple and strict: send a decided difference around the
machine's own encoding and demand that it come back the same difference.

Lay the encoding out plainly. The machine writes its numbers in a code --- at bottom a \emph{binary} code, marks over a
two-symbol alphabet --- and a code is trustworthy only if it is uniquely decodable, if no codeword is confused with the
start of another. Kraft~\cite{kraft1949} gave the exact condition under which such a code exists, a finite inequality on
the codeword lengths; and the machine's code satisfies it by construction, so that every value it writes down it can
read back without ambiguity. Encoding, in other words, is not allowed to lose a distinction. A difference the machine
tanged before encoding must be a difference it can tange after decoding; the round-trip is required to be faithful, and
a code that failed this test would be no representation at all but a leak.

Now close the loop and watch what is demanded. The machine takes a value --- a finite condition, a bracket --- encodes
it into binary, and decodes it back. Repeatability is the requirement that what returns is numerically the same: the
same bracket, the same decided differences, nothing added and nothing lost in the passage. When a value survives this
round-trip unchanged, it has earned the name the machine gives it; it is \emph{representable}, a value the machine can
hold, ship, and recover without drift. The first loop is thus the machine's first act of self-consistency: it does not
merely build numbers, it checks that its own representation does not corrupt them. And a value that came back altered
would not be a failure to be smoothed over --- it would be a signal, the first faint form of the residue the next
chapter is entirely about.

Notice the shape of the test, because the book will run it again at every scale. Build something; send it around a
transformation; ask what came back changed. Here almost nothing changes --- the encoding is faithful, the loop closes
clean, the number returns itself --- and that clean closure is exactly what lets the machine trust its own numbers. But
the interesting cases, the ones the whole instrument is built to read, are the loops that \emph{do not} close clean:
where something survives the round-trip altered, and the alteration is the measurement. The first loop is the training
case, the one where the answer is ``nothing,'' so that when a later loop returns ``something,'' the machine knows a
reading when it sees one.

In the two verbs the first loop is a tange that certifies a funge. The round-trip tanges the returned value against the
value sent --- decides whether they are the same --- and when the decision comes back \emph{same}, the machine funges
them: the original and its round-tripped image are pooled as one number, two representations of a single value. The
loop is faithful exactly when this funge closes with nothing left over. The machine can now tell apart, count, and
represent --- and it has one loop's worth of evidence that its representation is honest. What it does not yet have is a
loop that leaves a remainder. That remainder is the antagonist of the whole book, and the next chapter is where it
first refuses to disappear.
